% SOURCE_AUTHORITY: sga5_fr_workpass.tex, lines 9012--9042 (LF-normalized unit).
% SOURCE_SHA256: 791F4EFFC5E02832D5D77ED03518C8156D6F07E4C8238B03545DB93D883FBB28
\subsubsection*{3.3}
Por passagem ao limite, é evidente que as classes de cohomologia \(\ell\)-ádica do número anterior verificam as mesmas propriedades formais que as classes de cohomologia estudadas na Exposição IV, a saber:

\begin{enumerate}[label=(\roman*)]
\item Se \(f:X\to Y\) é um morfismo próprio de pré-esquemas suaves sobre \(k\) e \(Z\) um ciclo sobre \(X\), temos:
\[
f_*\bigl(\gamma_X(Z)\bigr)=\gamma_Y(f_*(Z)).
\tag{*}
\]

\item Se \(f:X\to Y\) é um morfismo de pré-esquemas suaves sobre \(k\) e \(Z\) um ciclo sobre \(Y\), temos:
\[
\gamma_X(f^*Z)=f^*\bigl(\gamma_Y(Z)\bigr).
\tag{**}
\]

\item Se \(X\) é um pré-esquema suave sobre \(k\) e \(Z_1\) e \(Z_2\) são dois ciclos sobre \(X\), temos:
\[
\gamma_X(Z_1\cdot Z_2)=\gamma_X(Z_1)\cdot \gamma_X(Z_2),
\tag{**}
\]
onde o ponto do primeiro membro denota o produto de intersecção dos ciclos e o ponto do segundo membro o produto cup em cohomologia.

\item Dois ciclos algebricamente equivalentes sobre um \(k\)-pré-esquema liso têm a mesma classe de cohomologia \(\ell\)-ádica.
\end{enumerate}

\begin{flushleft}
\(\quad (*)\) Esta compatibilidade não é demonstrada em (SGA \(4^{1/2}\) Cycle), mas não é difícil deduzi-la dos resultados de (loc. cit.).

\(\quad (**)\) No caso sem excesso, veja-se (SGA \(4^{1/2}\) Cycle 2.3.8) para um enunciado preciso.
\end{flushleft}
